%% Per-system sensitivity of the carrier search.  Table fragment generated by
%% persys_r17.py; the caption quotes only macros other generators publish.
\section{Per-system sensitivity}
\label{app:persys}

Table~\ref{tab:persys} gives the sensitivity reached toward each system of the
carrier search together with the three things that decide what that number
means: where in frequency it was obtained, over how wide a channel, and
whether the archive holds a second epoch capable of confirming a crossing
there. The limit quoted for a system is that of its best window, so a system
whose coverage is wide is not penalised for the shallowest window in it; the
whole distribution, window by window, is the per-window result the sensitivity
of \S\ref{sec:pninetysel} is drawn from.

Three features of the table bear on how it should be read. The limits span
more than two orders of magnitude, so a single survey depth would describe no
part of the sample. The frequency at which a system's best limit falls is set
by what the archive happens to hold toward it and not by any property of the
star. And a limit in the first column says nothing about whether a crossing at
that system could have been followed up: the last column does, and for
\RcNoRepSys{} of the \NSysClassA{} systems no sensitivity would have been
enough, because no second block covers a frequency the first one did.

\begin{table*}
\centering
\caption{\textbf{Per-system sensitivity of the carrier search.} For each of
the \NSysClassA{} systems the fine-channel search covers, the
$\mathrm{EIRP}_{90}$ of its best window --- the equivalent isotropic radiated
power at which nine injected carriers in ten are recovered, $\EirpNinetyMultA
\,P_{\rm trig}$ with that window's own annual-parallax and smearing retention
applied --- with that window's centre frequency $\nu$ and native channel width
$\Delta\nu$, and the system's merged frequency coverage $\Delta\nu_{\rm cov}$.
Rows are ordered by the limit, which spans
$\EirpNinetySysLoA$--$\EirpNinetySysHiA$\,W; the median system is covered over
\SysUnionAMed\,GHz, and channel widths across the search run from \ChanLoA{}
to \ChanHiA\,kHz. The last column records whether two execution blocks cover a
common frequency more than one day apart (d) or more than a year apart (yr):
\RcConfSysDay{} systems reach the first and \RcConfSysYr{} the second, and
these are lower bounds, since two blocks can be shown to be separated but
never shown to be simultaneous. For the remaining \RcNoRepSys{} systems a
crossing could not have been confirmed at any sensitivity. A gravitationally
bound pair is one system and takes its primary's name.}
\label{tab:persys}
\footnotesize
%% GENERATED by persys_r17.py -- do not hand-edit.
%% One row per system of the carrier search, two systems per printed line,
%% ordered by the limit.  Every value is computed from the per-window results.
\begin{tabular}{@{}l@{\hspace{5pt}}r@{\hspace{5pt}}r@{\hspace{5pt}}r@{\hspace{5pt}}r@{\hspace{5pt}}c@{\hspace{12pt}}l@{\hspace{5pt}}r@{\hspace{5pt}}r@{\hspace{5pt}}r@{\hspace{5pt}}r@{\hspace{5pt}}c@{}}
\hline
System & $\mathrm{EIRP}_{90}$ & $\nu$ & $\Delta\nu$ & $\Delta\nu_{\rm cov}$ & conf. & System & $\mathrm{EIRP}_{90}$ & $\nu$ & $\Delta\nu$ & $\Delta\nu_{\rm cov}$ & conf. \\
 & ($10^{14}$\,W) & (GHz) & (kHz) & (GHz) &  &  & ($10^{14}$\,W) & (GHz) & (kHz) & (GHz) &  \\
\hline
LP 876-10 & 0.87 & 345.83 & 488 & 1.7 & d & HD 170773 & 17 & 230.13 & 488 & 7.4 & d \\
$\beta$~Pic & 1.0 & 230.52 & 15.3 & 4.4 & yr & HD 61005 & 19 & 230.52 & 244 & 19 & d \\
AU Mic & 1.2 & 230.56 & 488 & 3.5 & yr & HR 1010 & 19 & 230.53 & 977 & 2.8 & d \\
$\tau$~Cet & 1.2 & 230.55 & 488 & 1.8 & d & LHS 1140 & 19 & 230.00 & 488 & 1.8 & d \\
HD 69830 & 1.2 & 230.06 & 488 & 3.8 & d & $\kappa$~CrB & 19 & 230.57 & 977 & 1.7 & -- \\
61 Vir & 1.3 & 345.79 & 488 & 1.8 & yr & HD 155555C & 21 & 346.01 & 488 & 1.7 & -- \\
$\epsilon$~Eri & 1.4 & 229.98 & 977 & 2.3 & yr & L 836-122 & 22 & 346.00 & 488 & 1.7 & -- \\
$\eta$~Crv & 1.4 & 345.80 & 61.0 & 7.2 & yr & CD$-$57 1054 & 23 & 345.99 & 488 & 1.7 & -- \\
TRAPPIST-1 & 1.6 & 346.58 & 488 & 1.7 & d & CP$-$72 2713 & 24 & 345.12 & 488 & 5.8 & -- \\
AT Mic AB & 1.7 & 346.04 & 488 & 1.7 & -- & HD 14055 & 24 & 230.11 & 488 & 7.2 & d \\
GJ 3379 & 2.2 & 229.96 & 488 & 1.9 & d & TWA 7 & 24 & 345.99 & 244 & 0.87 & d \\
GJ 14 & 3.1 & 227.22 & 488 & 3.5 & d & GJ 2006A & 24 & 346.02 & 488 & 1.7 & -- \\
HN Lib & 3.9 & 229.99 & 488 & 1.8 & d & Wolf 219 & 26 & 230.13 & 488 & 1.7 & d \\
HD 10647 & 4.4 & 345.79 & 488 & 6.1 & yr & LP 736-15 & 28 & 345.82 & 977 & 1.7 & -- \\
HD 202628 & 5.4 & 230.51 & 488 & 1.8 & yr & TYC 2211-1309-1 & 28 & 346.00 & 488 & 1.7 & -- \\
HD 53143 & 5.4 & 230.53 & 977 & 1.8 & yr & UCAC3 176-23654 & 29 & 346.02 & 488 & 1.7 & -- \\
HD 172555 & 5.8 & 230.53 & 15.3 & 5.3 & -- & PM J03433+1958 & 31 & 230.13 & 488 & 1.7 & -- \\
HD 207129 & 6.4 & 230.57 & 488 & 1.8 & yr & UCAC2 20312880 & 32 & 345.96 & 488 & 1.7 & -- \\
GJ 849 & 6.8 & 230.02 & 488 & 1.9 & d & LP 353-51 & 33 & 346.02 & 488 & 1.7 & -- \\
TWA 3A & 7.6 & 230.54 & 122 & 0.44 & d & Barta 161 12 & 34 & 346.02 & 488 & 1.7 & -- \\
HD 139664 & 7.8 & 230.55 & 488 & 1.7 & -- & 51 Eri & 37 & 230.53 & 1953 & 1.7 & -- \\
HD 92945 & 8.3 & 355.11 & 488 & 3.7 & d & HD 23484 & 38 & 230.53 & 977 & 1.8 & d \\
TX PsA & 8.5 & 346.03 & 488 & 1.7 & -- & HD 139084 & 54 & 346.01 & 488 & 1.7 & -- \\
HD 285968 & 9.1 & 229.97 & 488 & 1.8 & d & WD 0407+179 & 84 & 230.13 & 488 & 1.7 & -- \\
HD 161868 & 9.9 & 227.22 & 488 & 7.4 & d & HD 48370 & 89 & 230.73 & 122 & 3.3 & yr \\
HD 38858 & 13 & 230.50 & 488 & 1.8 & yr & HD 45184 & 93 & 230.53 & 977 & 1.7 & d \\
HD 110411 & 13 & 219.58 & 244 & 4.7 & yr & ALMA J153702653-33192492 & 99 & 345.98 & 488 & 14 & d \\
HIP 490 & 13 & 230.53 & 244 & 0.86 & -- & HD 31392 & 130 & 230.52 & 977 & 1.7 & d \\
HIP 10679 & 14 & 230.55 & 15.3 & 0.11 & -- & HD 89452 & 260 & 230.53 & 977 & 1.7 & d \\
HD 107146 & 14 & 355.14 & 488 & 3.7 & d & HD 377 & 370 & 230.54 & 977 & 1.7 & d \\
LP 476-207 & 15 & 346.02 & 488 & 1.7 & -- &  &  &  &  &  &  \\
\hline
\end{tabular}

\end{table*}
